\chapter{Methodology and Work Plan Overview}
\setcounter{section}{0}

This chapter describes how the project was actually carried out: which materials and resources were
used, which technical methods are applied where, what process governed the day-to-day decisions, in
what sequence the work was done, and which constraints shaped every choice. The substance of
\emph{what was built} is reserved for the Technical Development chapter
(Section~\ref{sec:system_design} and following), and the substance of \emph{what was discovered
experimentally} is reserved for the Experiments chapter (Section~\ref{sec:exp_methodology} and
following); this chapter explains the \emph{way of working} that connects them.

% =============================================================================
\section{Materials and Resources}
\label{sec:plan_materials}
% =============================================================================

The project required four classes of resources: hardware, software stack, datasets, and
supervision. Each class is described briefly below and cross-referenced to the chapters where the
technical detail lives.

\paragraph{Hardware.} All training, evaluation, and landmark extraction were performed on a single local workstation
whose specification is the operational environment of Section~\ref{sec:operational_environment};
the full table of components and version constraints is Table~\ref{tab:operational_environment} and
is not repeated here. The relevant element for the work-plan discussion is the GPU: an NVIDIA
RTX~3050~Ti Laptop with 4\,GB of VRAM. The decision to keep the entire training pipeline on a local
GPU rather than on a cloud platform such as Google Colab or AWS was deliberate. Local execution
gives bit-for-bit reproducibility between consecutive runs, costs nothing per run (which mattered
because the iterative methodology described in Section~\ref{sec:plan_methodology} relaunches
experiments often), and removes the network-dependent latency that would otherwise penalise short
interactive runs. The 4\,GB VRAM ceiling is, in turn, the dominant hardware constraint discussed in
Section~\ref{sec:plan_constraints}.

\paragraph{Software stack.} The training and evaluation pipeline runs on Python~3.10 with PyTorch~2.x and CUDA~12.x. Pose
estimation uses MediaPipe Hands~0.10+, video decoding uses OpenCV~4.8+, and version control is
provided by Git with GitHub as remote host. The mobile client is built with React Native and Expo;
the inference server uses FastAPI~0.111 with Uvicorn~0.29. Documentation is authored in \LaTeX. The
exact version constraints are locked in \texttt{pyproject.toml} for the training package and in
\texttt{04\_server/requirements.txt} for the inference server, so that the operational environment
is reproducible from a clean clone of the repository (cross-reference:
Section~\ref{sec:operational_environment} and the reproduction guide at
\texttt{docs/reproducibility\_guide.md}).

\paragraph{Datasets.} Six datasets are used across the project: EgoHands and the Kaggle Hand Dataset for hand
localisation, Sign Language MNIST and the ASL Alphabet for letter classification, and WLASL and
MS-ASL for word recognition. Each dataset is formally defined in Section~\ref{sec:datasets}, where
attribute schema, licence considerations, and the role of each dataset in the experimental pipeline
are described. They are listed here only to make the resource inventory complete; the work-plan
discussion treats the datasets as external inputs whose acquisition, audit, and preparation costs
are accounted for in the corresponding phases of Section~\ref{sec:plan_workplan}.

\paragraph{Supervision and consultation resources.} The project was supervised by Mois\'es Mart\'inez Mu\~n\'oz through scheduled tutor meetings held
weekly or biweekly depending on the phase of the work. The most consequential technical decisions
of the project such as the pivot from pixel-based letter recognition to landmark-based word
recognition, the adoption of stratified $K$-fold cross-validation in place of test-time
augmentation, the cross-dataset extension to MS-ASL50 in addition to WLASL, and the introduction of
the multi-seed protocol, did not originate as solo decisions but as outcomes of these meetings. The
archived meeting notes preserve the directives that triggered each shift, and an internal audit
document was used to trace the meeting references cited in this chapter. Consultation resources
beyond the tutor consist of the literature reviewed in chapter~\ref{ch:sota} (notably the WLASL
benchmark, MediaPipe Hands, and the Bahdanau attention formulation) and a small number of
open-source repositories used as reference implementations: the WLASL benchmark distribution and
the MediaPipe Hands toolkit.

% =============================================================================
\section{Methods}
\label{sec:plan_methods}
% =============================================================================

This section is a compact catalogue of the technical methods the project relies on, presented
without re-justification. The theoretical background of each method is the subject of
chapter~\ref{ch:sota}, and the empirical evidence for each one is the subject of the Experiments
chapter (Section~\ref{sec:exp_methodology} and following); the role of this section is to make the
inventory explicit so that the rest of the chapter can refer to it.

The perception frontend uses MediaPipe Hands, a top-down two-stage pipeline that first detects the
palm and then regresses 21 three-dimensional keypoints per detected hand. The choice to operate on
landmarks instead of raw pixels is the central methodological decision of the project and is
justified in detail in Section~\ref{sec:exp_phase1}; the relevant point here is that all downstream
methods consume a landmark vector of 63 features per frame.

The temporal model is a bidirectional Gated Recurrent Unit with 256 hidden units per direction and
two stacked layers, followed by a Temporal Attention head implemented as an additive
(Bahdanau-style) softmax over the per-frame BiGRU outputs. A linear classifier projects the
attention-weighted context vector to the gloss vocabulary. This architecture is referred to
throughout the document as the GRU+Attention model and totals approximately 2.26\,million
parameters at the configurations reported in Section~\ref{sec:exp_methodology}. Two alternative
architectures appear in the comparison tables of Section~\ref{sec:exp_phase2}: a plain
bidirectional GRU baseline and a Transformer encoder with four layers and eight heads. They share
the input pipeline and the loss function with the champion model and differ only in the
sequence-modelling block.

Optimisation is performed with AdamW at learning rate $1\times10^{-3}$, weight decay
$5\times10^{-4}$, and label smoothing $0.1$ on the cross-entropy loss. The learning-rate schedule
is CosineAnnealingLR. Early stopping monitors the validation Top-1 accuracy with patience 30 epochs
and a maximum run length of 150 epochs. Skeletal augmentation during training combines rotation,
scaling, translation, additive Gaussian noise, and horizontal flip, applied per clip on the
landmark coordinates rather than on pixels.

Evaluation uses two complementary protocols. Stratified 5-Fold cross-validation runs on the WLASL30
vocabulary, with the dataset rebuilt by filtering WLASL100 down to the same 30 classes (332
samples) and the folds drawn so that the per-class distribution is preserved. The multi-seed
protocol runs the same configuration with three random seeds (42, 123, 2026) and reports the Top-1,
Top-5, Macro-F1, and Balanced Accuracy metrics as mean$\,\pm\,$standard deviation. Test-time
augmentation was deliberately excluded from the protocol after supervisory review identified it as
a source of biased inflation; the rationale is recorded in Section~\ref{sec:exp_methodology}.

Reproducibility is a method as much as a property. Every training and evaluation run produces an
\texttt{experiment\_manifest.json} that captures the git commit, the resolved configuration, the
hardware identifier, the wall-clock time, and the final metrics. Hyperparameters and dataset paths
live in YAML files under \texttt{configs/}, never in code, with a base configuration defining the
defaults shared across experiments and per-experiment files overriding only the fields they need.
The full reproduction procedure is documented in \texttt{docs/reproducibility\_guide.md}.

The exploratory phase of the project relied on two methods that no longer carry results in the
Experiments chapter but are listed here for completeness. The first is YOLOv8n hand detection, used
during Phase~B (see Section~\ref{sec:plan_workplan}) to localise hands from raw RGB frames. The
second is CNN-based pixel classification of static ASL letters. Both methods reached working
accuracy on their respective tasks and both were deprecated when the project pivoted to
landmark-based word recognition; Section~\ref{sec:exp_phase1} explains why.

% =============================================================================
\section{Methodology}
\label{sec:plan_methodology}
% =============================================================================

The project follows an iterative, experiment-driven methodology rather than a fixed waterfall plan,
because the technical direction shifted substantially after the early exploratory phase and any
plan that did not allow for that shift would have been discarded in February. The methodology is
not borrowed from CRISP-DM, Scrum, or any other process-software framework: those frameworks were
considered, but the development log documents a different and simpler pattern, applied repeatedly
across the lifetime of the project. That pattern is the basis of the methodology described here.

Each meaningful step of the project is organised as a four-stage loop. The loop starts with a
\emph{hypothesis} that is verifiable in principle: for example, ``landmark-based representation
outperforms pixel-based representation on WLASL30 within the available data budget,'' or ``a
sequence length of 64 frames produces lower seed-to-seed variance than 32 or 96.'' The loop
continues with \emph{implementation}: the relevant code changes are made in the \texttt{src/slr/}
package, augmented with whichever new YAML configuration the experiment requires, and committed to
git so that the run is anchored to a specific repository state. The loop then runs the
\emph{experiment}: the training script is launched, the manifest is written, and the metrics are
persisted alongside the model checkpoints in the experiment artifact directory. Finally, the loop
closes with \emph{reflection}: a devlog entry records what was attempted, what the result was,
whether the hypothesis was confirmed, and what the next step should be. If the hypothesis fails,
the next iteration starts from a new hypothesis informed by the failure rather than from the
original plan.

The pivot recorded in Section~\ref{sec:exp_phase1} is the most consequential application of this
loop. The project began with a pixel-based pipeline that targeted hand detection with YOLOv8n
followed by per-frame letter classification. Hand detection reached approximately 98\,\% mAP@0.5 on
the joint EgoHands+Kaggle test split and was, by any reasonable measure, working. Letter
classification was equally tractable on the static-frame benchmarks. The combined pipeline,
however, surfaced a structural observation that no incremental tuning would have removed: the
bottleneck of the system was not localising the hand but representing the dynamic content of the
sign. The four-stage loop responded to that observation by retiring the pixel-based recognition
stage and replacing it with the landmark-based word-recognition pipeline that is the subject of
every subsequent chapter. The pivot was a methodological consequence, not an arbitrary scope
change: the iterative loop is what made the pivot possible, and the development log makes the
rationale auditable in retrospect.

Two further properties of the methodology are worth making explicit because they show up repeatedly
in the Technical Development and Experiments chapters that follow. The first is the role of tutor
meetings as control points. The most impactful technical decisions of the project were made jointly
with the supervisor at scheduled meetings rather than unilaterally between iterations. These
decisions included removing test-time augmentation from the evaluation protocol after supervisory
review, restructuring the State of the Art chapter as a hierarchical
CNN$\rightarrow$RNN$\rightarrow$LSTM$\rightarrow$GRU theory progression, adopting Google's AI
Principles during the ethics review phase, and splitting the experimental work into quantitative
and qualitative tracks during the final experimental review. The meetings act as low-frequency
external feedback that the four-stage loop, on its own, would not provide.

The second property is the introduction of the multi-seed protocol described in
Section~\ref{sec:plan_methods}. The protocol is a direct response to a tutor concern about
reproducibility raised in early April~2026: a single-seed Top-1 figure on the 34-sample WLASL30
test split is bounded below in resolution by the granularity of one misclassification
($\approx$2.94\,pp), and a single number at that resolution is indistinguishable from
initialisation noise. The methodology absorbed the concern by adding a new ``aggregate over three
seeds'' step to every headline experiment going forward and by retroactively applying it to the
sequence-length sweep and to MS-ASL50. This is a concrete example of the iterative methodology
incorporating supervision: the protocol that distinguishes a defensible thesis number from a noisy
point was added in response to a single conversation, propagated through the codebase by adding a
small number of YAML files, and reflected in Section~\ref{sec:exp_methodology} as the statistical
backbone of the experimental claims.

Documentation is treated as a continuous side effect of the loop rather than as a separate phase.
The development log contains one entry per meaningful session, dated, scoped to a specific session
name, and indexed centrally. The log is the primary source from which the work plan in
Section~\ref{sec:plan_workplan} is reconstructed; without it, the chronology of the project would
be unrecoverable from the code alone.

% =============================================================================
\section{Work Plan}
\label{sec:plan_workplan}
% =============================================================================

The executed work plan was reconstructed from the original proposal, recorded development notes,
experiment artifacts, and supervision notes. The plan is organised into seven phases (A--G). The
first two phases correspond closely to the proposal-era plan, while the remaining phases reflect
the methodological pivot from static, pixel-based recognition to landmark-based word-level ASL
recognition.

\paragraph{Phase A: Project framing and proposal.} The project formally started in the first week of January~2026 with problem definition, preliminary
literature review, scope selection, and proposal preparation. This phase defined the accessibility
motivation, the initial system architecture, the first work schedule, and the research direction.
No model training was carried out at this stage; the main output was the approved project proposal
and its initial planning material.

\paragraph{Phase B: Hand detection and static letter recognition.} The initial technical work explored the proposal's first recognition route: localising hands in
images and classifying static ASL letters. A YOLOv8n detector was trained on a combined EgoHands
and Kaggle Hand Dataset corpus, and CNN-based classifiers were explored on static ASL alphabet
datasets. These components reached working performance on their own tasks, but they also revealed a
limitation in the original direction: the central research challenge was not hand localisation, but
representing the temporal structure of complete signs.

\paragraph{Phase C: Landmark pivot and WLASL30 baseline.} In late January and February 2026, the project pivoted from pixel-based letter recognition to
word-level recognition using MediaPipe hand landmarks. WLASL30 was built as the first controlled
benchmark, and several sequence models were compared under the same landmark input representation.
The GRU+Attention architecture emerged as the strongest baseline because it combined a recurrent
temporal model with an attention mechanism able to focus on the most informative frames. After
supervisory review, test-time augmentation was excluded from the final evaluation protocol, and the
WLASL30 result was strengthened through stratified $K$-fold validation.

\paragraph{Phase D: Vocabulary scaling and ethical framework.} Once the GRU+Attention model had been selected, the next phase tested how the approach behaved as
the WLASL vocabulary increased beyond the initial 30-class setting. The model was retrained on
progressively larger WLASL subsets to quantify the effect of vocabulary size and per-class data
scarcity. In parallel, the ethical and social-impact chapter was developed, including the
accessibility motivation, privacy implications of landmark-based representations, bias
considerations, reproducibility, and the responsible-deployment limits of the prototype.

\paragraph{Phase E: Runtime prototype.} In April 2026, the technical work expanded from offline training to a runtime proof of concept. A
mobile client and FastAPI server were integrated so that short clips could be captured, sent for
inference, and returned as Top-$K$ word predictions. Additional prototype features included
text-to-speech, model selection from the client, and a heuristic continuous-prediction demo mode
based on overlapping windows. This phase demonstrates end-to-end feasibility, but it is not treated
as the main research result; the quantitative contribution remains the landmark-based recognition
pipeline and its validation.

\paragraph{Phase F: Refactoring and cross-dataset validation.} The notebook-based experimental code was consolidated into a modular Python package, with
configuration-driven training, structured experiment manifests, reusable evaluation scripts, and
preserved checkpoints and plots. This refactor improved reproducibility and made it possible to
rerun the same pipeline across datasets. MS-ASL50 was then added as an independent validation
benchmark, and the final evaluation protocol incorporated multi-seed runs, a sequence-length sweep,
cross-dataset validation, and comparison against the WLASL results.

\paragraph{Phase G: Documentation and qualitative validation design.} The final phase focused on thesis writing, result consolidation, figure and table preparation, and
the design of a qualitative validation protocol for the runtime prototype. The experimental chapter
was organised into quantitative and qualitative tracks: the quantitative track reports completed
model evaluations, while the qualitative track defines the participant protocol and analysis plan
for future usability evidence. This phase remains open at the time of this intermediate submission
and will be completed with the final thesis revision and qualitative data integration.

Figure~\ref{fig:work_plan_timeline} summarises the executed work plan. The chart follows the seven
phases described above and highlights the main deviation from the initial proposal: the pivot from
static letter recognition to landmark-based word-level recognition.

\image{plan/work_plan_gantt.png}{0.95}{Project Gantt chart showing the executed work plan grouped into seven phases, from project framing and initial hand-detection experiments to landmark-based word recognition, cross-dataset validation, prototype integration, and final documentation. Source: self-created with GanttProject.}{fig:work_plan_timeline}

\FloatBarrier

% =============================================================================
\section{Constraints and Limitations}
\label{sec:plan_constraints}
% =============================================================================

This section catalogues the constraints that shaped the work plan and the methodology rather than
the experimental results. Limitations of the experimental conclusions themselves, including
standard deviations, test-set granularity, and single-architecture cross-dataset validation, are
treated separately in Section~\ref{sec:exp_limitations}.

The first constraint is time. The project was developed in parallel with a software-engineering
internship, which restricted the number of contiguous hours that could be dedicated to it in any
given week. The iterative methodology described in Section~\ref{sec:plan_methodology} is in part a
response to this constraint: a four-stage loop with an explicit devlog entry per session is what
makes it possible to resume the project after a multi-day pause without losing the local context.
Without the log, the cost of re-entering a problem after an interruption would have dominated the
cost of solving it.

The second constraint is hardware. Every training and evaluation run executed on a single NVIDIA
RTX~3050~Ti Laptop with 4\,GB of VRAM. This ceiling propagated into multiple design choices:
\texttt{max\_frames}\,=\,64 was selected partly because longer sequences exceeded the memory budget
at batch size 32, the GRU+Attention model was sized at approximately 2.26\,million parameters
rather than the parameter counts typical of large Transformers, and the Transformer baseline that
does appear in Section~\ref{sec:exp_phase2} uses a relatively narrow $d_{\mathrm{model}}=256$ for
the same reason. A cloud GPU was considered and rejected: per-run cost predictability and
bit-for-bit reproducibility between consecutive runs were judged more valuable than the additional
VRAM.

The third constraint is data availability. Both word-recognition datasets used in the project are
reconstructed from external sources rather than distributed directly, and both have lost a
measurable fraction of their original clips. WLASL relies on URLs that have decayed since the
dataset was published, and the MediaPipe detection rate on the surviving WLASL clips is
approximately 69\,\%, which limits the effective per-clip information content. MS-ASL is
reconstructed from per-instance YouTube URLs, and the reconstruction audit records a download
success rate of 65.9\,\% (16{,}816 of 25{,}513 clips). The MS-ASL50 subset, rather than the more
conventional MS-ASL30, is a direct response to this constraint: only six of the top-30 MS-ASL
classes meet the strict 30/5/5 train/val/test threshold from current YouTube state, so an MS-ASL30
reproduction would have been statistically indefensible (cf.\ Section~\ref{sec:exp_msasl}).

The fourth constraint is evaluator availability. Until the qualitative validation protocol
described in Section~\ref{sec:exp_qualitative} executes, the author is the only evaluator and the
only experimenter on the project. This is a standard configuration for an undergraduate
final-degree project and is acceptable for the quantitative work, where the evaluation is fully
automated and reproducible from a manifest. It does, however, mean that the cognitive bias of the
experimenter is uncontrolled at the design level. The qualitative protocol is structured so that
the analysis plan is fixed in advance of data collection, in the spirit of pre-registration,
precisely to mitigate this bias once the protocol runs.

The fifth constraint is linguistic scope. The project validates a single sign language, American
Sign Language, and uses ASL exclusively for both the WLASL and MS-ASL benchmarks. The pipeline
contains no ASL-specific assumptions, and every component that consumes landmarks could in
principle be retrained on a different sign language given a comparable benchmark. The validation
has not been done, however, and any claim about transfer to Spanish Sign Language, British Sign
Language, or any other variety is therefore future work.

The sixth and most general constraint is that the iterative, experiment-driven methodology
described in Section~\ref{sec:plan_methodology} is, by construction, incompatible with the literal
execution of the Gantt that the project proposal recorded in January~2026. The Gantt scheduled hand
localisation, static handshape recognition, dynamic sign recognition, temporal fusion and decoding,
language generation, mobile prototype, component integration, and validation tests as a sequence of
rigid intervals; the methodology actually applied retired the Static Handshape Recognition track in
favour of word-level recognition, replaced the proposed PaliGemma language-generation step with a
direct gloss-to-text mapping, and folded the Temporal Fusion stage into the GRU+Attention model.
This divergence is not a planning failure but a deliberate property of the methodology. The tutor
meeting notes document the moments at which each scope change was approved, and the devlog and
meeting notes together act as the audit trail that links the proposal-era plan to the executed
plan.
